\section{The Bailout Protocol}
\label{sec:bailout-protocol}

\subsection{Overview and Purpose}

The Bailout Protocol is the mandatory human accountability mechanism embedded in every node of the BX3 recursive architecture. It ensures that when an autonomous agent encounters a condition it cannot resolve with certainty — whether anticipated or not — the exception does not vanish into silent failure or unchecked autonomous action. Instead, it propagates upward through the hierarchy, bypassing all Machine actors, until it reaches a Human Accountability Anchor. This structural property makes human accountability indestructible: it survives recursion, delegation, and network disconnection.

The protocol addresses the \emph{black box problem} in autonomous systems — the failure mode in which autonomous actions proceed without any human knowing, authorizing, or being able to account for them. The Bailout Protocol is not an error handler. It is a universal escalation path that applies to every condition a node cannot resolve with certainty, regardless of category, severity, or whether the condition was anticipated in the node's programming.

Every node — from the lowest-grade sensor to the Human Root at the apex of the hierarchy — carries an identical copy of the Bailout Protocol embedded in its Worksheet. The protocol cannot be removed, disabled, or overridden by any Machine actor. It is structurally indestructible.

\subsection{Formal Definitions}

\begin{definition}[Human Accountability Anchor (HAA)]
A node in the BX3 hierarchy that is \emph{legally authorized} to resolve a pending exception by issuing a Purpose directive. The Human Root is the sole terminal HAA in the architecture. All other nodes are Machine Accountability Anchors (MAAs).
\end{definition}

\begin{definition}[Machine Accountability Anchor (MAA)]
A node in the BX3 hierarchy that \emph{cannot} serve as a terminal accountability point for exception resolution, regardless of its operational authority. MAAs may evaluate and attempt to resolve triggers locally, but they cannot close a trigger that requires human judgment. Every non-Human Root node is an MAA by definition.
\end{definition}

\begin{definition}[Bailout Trigger]
A structured exception event generated by a node when it encounters a condition it cannot resolve with certainty. The Bailout Trigger carries a full Trigger Package and propagates asynchronously up the parent chain until it reaches an HAA.
\end{definition}

\begin{definition}[Trigger Package]
A self-contained data structure emitted by a node upon firing a Bailout Trigger. It contains: the originating node identifier, the triggering event, the applicable Trigger Condition classification, a confidence score, the full node state snapshot at time of trigger, a proposed resolution (if determinable), the propagation chain history, and the current protocol state.
\end{definition}

\begin{definition}[Safety Envelope]
A set of hard operational constraints (physical, regulatory, ethical) that the Bounds Engine is prohibited from permitting any action to violate. Violations, whether actual or predicted by sandbox simulation, constitute a Bailout Trigger.
\end{definition}

\begin{definition}[Confidence Score]
A real number $c \in [0.0, 1.0]$ that quantifies the originating node's certainty that the triggered condition requires Human Root resolution. A score of 0.0 indicates the node is propagating for safety only (uncertain). A score of 1.0 indicates the node is certain no machine resolution is possible.
\end{definition}

\subsection{The State Machine}

The Bailout Protocol operates as a deterministic finite state machine (DFSM) over six mutually exclusive states. The state machine is embedded identically in every node's Worksheet and governs how that node processes, propagates, and resolves bailout events.

\bigskip
\noindent\textbf{State Space:} $\{ \text{IDLE}, \text{BOUNDED}, \text{BAIL\_PENDING}, \text{ESCALATING}, \text{HUMAN\_ACKNOWLEDGED}, \text{RESOLVED} \}$

\bigskip
\noindent\textbf{State Definitions:}

\begin{itemize}
  \item[\textbf{IDLE}] The node is operating within normal parameters. No bailout condition is present. The node is actively monitoring for trigger conditions but is not processing any exception.
  \item[\textbf{BOUNDED}] The node has detected a condition that \emph{may} require escalation. The condition is being evaluated by the local Bounds Engine. No trigger has been fired. The node is in a holding state — it will not execute any action that would alter the system state until evaluation completes.
  \item[\textbf{BAIL\_PENDING}] A Bailout Trigger has been fired. The Trigger Package is being assembled and is in the process of being transmitted to the parent node. The node is \emph{not} awaiting a response from the parent — the transmission is asynchronous and confirmed.
  \item[\textbf{ESCALATING}] The Trigger Package has been delivered to the parent node. The parent is actively evaluating the trigger. The originating node is holding and awaiting a directive from above.
  \item[\textbf{HUMAN\_ACKNOWLEDGED}] The trigger chain has reached the Human Root. The Human Root has received the Trigger Package and is actively reviewing it. The originating node is holding and awaiting a Purpose directive.
  \item[\textbf{RESOLVED}] The trigger has been closed. A resolution directive has been received (either from a Machine node that could resolve it locally, or from the Human Root), and the node is executing the directive. Upon completion of the directive execution, the state machine transitions to IDLE.
\end{itemize}

\subsection{State Transitions}

The state machine transitions are governed by the tuple $(S, s_0, \Sigma, T, F)$ where:
\begin{itemize}
  \item $S = \{\text{IDLE}, \text{BOUNDED}, \text{BAIL\_PENDING}, \text{ESCALATING}, \text{HUMAN\_ACKNOWLEDGED}, \text{RESOLVED}\}$
  \item $s_0 = \text{IDLE}$ (initial state on node boot)
  \item $\Sigma$ is the input alphabet (trigger conditions, delivery confirmations, directives, timeouts)
  \item $T: S \times \Sigma \rightarrow S$ is the total deterministic transition function
  \item $F = \{\text{RESOLVED}\}$ (accepting / terminal-adjacent state; RESOLVED transitions to IDLE on directive execution complete)
\end{itemize}

\bigskip
\noindent The transition function $T$ is defined as follows:

\begin{enumerate}
  \item[\textbf{$T(\text{IDLE}, \text{EVALUATE}) = \text{BOUNDED}$}] When the node's monitoring subsystem detects a condition that \emph{may} exceed the node's resolution certainty threshold, the node transitions to BOUNDED and initiates Bounds Engine evaluation. The node does not fire a trigger at this stage — it first evaluates whether the condition is within its operational authority.

  \item[\textbf{$T(\text{BOUNDED}, \text{LOCAL\_RESOLUTION\_FOUND}) = \text{IDLE}$}] If the Bounds Engine determines that the condition is resolvable locally without human intervention and without violating the Safety Envelope, the evaluation completes with no trigger fired. The node returns to IDLE.

  \item[\textbf{$T(\text{BOUNDED}, \text{TRIGGER\_CONDITION\_MET}) = \text{BAIL\_PENDING}$}] If the Bounds Engine evaluation determines that one or more Trigger Conditions are satisfied (Section \ref{sec:trigger-conditions}), a Bailout Trigger is fired. The node enters BAIL\_PENDING and begins assembling and transmitting the Trigger Package to its parent.

  \item[\textbf{$T(\text{BAIL\_PENDING}, \text{PACKAGE\_TRANSMITTED}) = \text{ESCALATING}$}] Upon confirmed asynchronous delivery of the Trigger Package to the parent node, the originating node transitions to ESCALATING. Delivery confirmation is obtained via the Recursive Spawning uplink's confirmation protocol.

  \item[\textbf{$T(\text{ESCALATING}, \text{LOCAL\_RESOLUTION\_RECEIVED}) = \text{RESOLVED}$}] If the parent node can resolve the trigger locally (the parent is an MAA with sufficient authority and the resolution is within Safety Envelope), the parent issues a directive that propagates back down the tree. The originating node receives the directive and transitions to RESOLVED to execute it.

  \item[\textbf{$T(\text{ESCALATING}, \text{ESCALATED\_UPWARD}) = \text{ESCALATING}$}] If the parent node determines it cannot resolve the trigger and is itself an MAA, it propagates the Trigger Package upward and issues a delivery confirmation to the originating node. The originating node remains in ESCALATING, holding its state, awaiting a directive from higher in the chain.

  \item[\textbf{$T(\text{ESCALATING}, \text{HUMAN\_ROOT\_REACHED}) = \text{HUMAN\_ACKNOWLEDGED}$}] If the propagation chain reaches the Human Root (the sole HAA), the Human Root sends an acknowledgment downward. Upon receipt of this acknowledgment at the originating node, the state transitions to HUMAN\_ACKNOWLEDGED.

  \item[\textbf{$T(\text{HUMAN\_ACKNOWLEDGED}, \text{DIRECTIVE\_RECEIVED}) = \text{RESOLVED}$}] Upon receipt of a Purpose directive from the Human Root, the originating node transitions to RESOLVED to execute the directive.

  \item[\textbf{$T(\text{HUMAN\_ACKNOWLEDGED}, \text{ESCALATED\_UPWARD}) = \text{HUMAN\_ACKNOWLEDGED}$}] If the Human Root itself propagates the trigger to a higher human authority (e.g., legal counsel, executive), the originating node remains in HUMAN\_ACKNOWLEDGED holding for the directive. This is a no-op transition on the originating node.

  \item[\textbf{$T(\text{RESOLVED}, \text{DIRECTIVE\_EXECUTED}) = \text{IDLE}$}] Upon successful execution of the resolution directive, the node returns to IDLE. The trigger is closed; the complete lifecycle is logged to the Ledger.
\end{enumerate}

\subsection{Bailout Trigger Conditions}
\label{sec:trigger-conditions}

A Bailout Trigger fires when \emph{any one} of three Trigger Conditions is satisfied. These conditions are checked by the Bounds Engine at the BOUNDED evaluation stage. They are mutually exclusive in classification but not in application — multiple conditions may fire simultaneously from the same event.

\medskip
\noindent\textbf{Condition 1 — Capability Boundary (CB):}

The node encounters a condition — a sensor reading, computation result, execution request, or environmental state — that falls outside its \emph{provisioned operational range}. The node lacks the data, the model, the training, or the authority to process the condition deterministically. This includes:
\begin{itemize}
  \item A sensor reading that is outside the range of any model the node possesses
  \item A requested action that the node has no provisioned skill to execute
  \item A logical condition the node cannot evaluate with certainty (e.g., a novel cross-domain dependency)
  \item A communication failure that prevents the node from obtaining required oracle data
\end{itemize}

Formally: $\text{CB} \iff \exists d \in D_\text{observed}: d \notin \text{Range}(M_\text{node}) \land \neg \text{CanObtain}(d)$

\medskip
\noindent\textbf{Condition 2 — Safety Envelope Violation, Predicted (SEVP):}

The Bounds Engine's sandbox simulation (P6 Projection) indicates that a proposed or requested action would cause a Safety Envelope parameter to be exceeded if executed. The Bounds Engine is architecturally prohibited from executing any action that violates the Safety Envelope. However, it is equally prohibited from silently discarding the request. The resolution is to trigger upward: the Human Root must authorize a Safety Envelope override, or the request must be modified.

Formally: $\text{SEVP} \iff \exists a \in A_\text{proposed}: \text{Simulate}(a) \rightarrow \exists p \in P_\text{safety}: \text{Violates}(a, p)$

\medskip
\noindent\textbf{Condition 3 — Accountability Boundary (AB):}

The node encounters a decision point where the legal, regulatory, or ethical implications of the decision exceed the node's \emph{provisioned accountability authority}. The node has operational authority to perform the action, but lacks the legal standing or ethical warrant to authorize it independently. Examples include: an action requiring landowner authorization, a financial commitment exceeding the node's authority threshold, an action requiring regulatory permits, or an action with liability implications.

Formally: $\text{AB} \iff \exists a \in A_\text{proposed}: \text{RequiresAuth}(a, \text{HUMAN}) \land \neg \text{HasAuth}(\text{node}, a)$

\medskip
\noindent\textbf{Trigger Condition Evaluation Algorithm:} The following algorithm is executed by the Bounds Engine upon entering the BOUNDED state:

\begin{algorithm}
\caption{Bailout Trigger Condition Evaluation}
\begin{algorithmic}[1]
\Procedure{EvaluateTriggerConditions}{event $e$, node $n$}
  \State $triggers \gets \emptyset$
  \State $range \gets \text{GetOperationalRange}(n)$
  \State $safety\_envelope \gets \text{GetSafetyEnvelope}(n)$
  \State $accountability\_scope \gets \text{GetAccountabilityScope}(n)$

  \If{$\exists d \in e.\text{data}: d \notin range \land \neg \text{CanObtain}(n, d)$}
    \State $triggers.\text{add}(\text{CB}, 1.0)$ \Comment{CB is certain}
  \EndIf

  \State $sim\_result \gets \text{RunSandbox}(e.\text{proposed\_action}, safety\_envelope)$
  \If{$\text{sim\_result}.\text{violates\_envelope}$}
    \State $confidence \gets \text{ComputeConfidence}(sim\_result, e)$
    \State $triggers.\text{add}(\text{SEVP}, confidence)$
  \EndIf

  \If{$\text{RequiresHumanAuth}(e.\text{proposed\_action}) \land \neg \text{HasAuth}(n, e.\text{proposed\_action})$}
    \State $triggers.\text{add}(\text{AB}, 0.9)$ \Comment{AB is near-certain}
  \EndIf

  \If{$triggers.\text{size} > 0$}
    \State \Return $($\textbf{TRIGGER\_CONDITION\_MET}$, triggers)$
  \Else
    \State \Return $($\textbf{LOCAL\_RESOLUTION\_FOUND}$, \emptyset)$
  \EndIf
\EndProcedure
\end{algorithmic}
\end{algorithm}

\subsection{Escalation Algorithm}

When a Bailout Trigger is fired, the following escalation algorithm governs the asynchronous propagation of the Trigger Package through the hierarchy.

\begin{algorithm}
\caption{Bailout Escalation Algorithm — Node-Level}
\begin{algorithmic}[1]
\Procedure{BailoutEscalate}{$node$ $n$, TriggerPackage $pkg$}
  \State $pkg.\text{propagation\_chain}.\text{append}(\{ \text{node}: n.\text{id}, \text{action}: \text{TRIGGER\_FIRED}, \text{ts}: \text{NowNS} \})$
  \State $parent \gets n.\text{parent}$
  \State $pkg.\text{status} \gets \text{BAIL\_PENDING}$

  \State \Comment{Asynchronous transmission to parent; do not block}
  \State $\text{TransmitAsync}(parent, pkg)$
  \State $pkg.\text{status} \gets \text{ESCALATING}$
  \State $n.\text{state} \gets \text{ESCALATING}$

  \State \Comment{Holding state — await directive from above}
  \State \textbf{await} $\text{Directive}(n)$
\EndProcedure

\Procedure{ParentEvaluate}{$node$ $parent$, TriggerPackage $pkg$}
  \State $pkg.\text{propagation\_chain}.\text{append}(\{ \text{node}: parent.\text{id}, \text{action}: \text{PARENT\_EVALUATING}, \text{ts}: \text{NowNS} \})$

  \If{$\text{CanResolveLocally}(parent, pkg)$}
    \State $directive \gets \text{GenerateDirective}(parent, pkg)$
    \State $\text{SendDirectiveDown}(pkg.\text{originating\_node}, directive)$
    \State $pkg.\text{status} \gets \text{RESOLVED}$
  \ElsIf{$\text{IsHAA}(parent)$} \Comment{parent is Human Root}
    \State $\text{SignalHumanRoot}(parent, pkg)$
    \State $pkg.\text{status} \gets \text{HUMAN\_ACKNOWLEDGED}$
    \State \textbf{await} $\text{PurposeDirective}(parent, pkg)$
  \Else$} \Comment{parent is MAA, cannot resolve}
    \State $\text{PropagateUpward}(parent.\text{parent}, pkg)$
  \EndIf
\EndProcedure
\end{algorithmic}
\end{algorithm}

\begin{algorithm}
\caption{Bailout Escalation Algorithm — Human Root-Level}
\begin{algorithmic}[1]
\Procedure{HumanRootProcess}{$node$ $hr$, TriggerPackage $pkg$}
  \State $\text{SurfaceAlert}(hr.\text{dashboard}, pkg)$
  \State $hr.\text{state} \gets \text{HUMAN\_ACKNOWLEDGED}$
  \State \textbf{await} $hr.\text{human\_decision}$

  \If{$hr.\text{decision} = \text{APPROVE\_WITH\_MODIFICATION}$}
    \State $directive \gets \text{BuildModifiedDirective}(pkg, hr.\text{modifications})$
    \State $\text{SendDirectiveDown}(pkg.\text{originating\_node}, directive)$
  \ElsIf{$hr.\text{decision} = \text{APPROVE\_AS\_PROPOSED}$}
    \State $\text{SendDirectiveDown}(pkg.\text{originating\_node}, pkg.\text{proposed\_resolution})$
  \ElsIf{$hr.\text{decision} = \text{DENY}$}
    \State $directive \gets \text{BuildNullDirective}(pkg)$
    \State $\text{SendDirectiveDown}(pkg.\text{originating\_node}, directive)$
  \ElsIf{$hr.\text{decision} = \text{ESCALATE\_TO\_HUMAN}$}
    \State $\text{PropagateToHigherHuman}(hr, pkg)$
  \EndIf

  \State $pkg.\text{status} \gets \text{RESOLVED}$
  \State $\text{LogToLedger}(pkg)$
\EndProcedure
\end{algorithmic}
\end{algorithm}

The escalation chain continues until the trigger reaches an HAA. Each MAA in the chain performs a local evaluation: if it can resolve, it issues a directive and the chain terminates downward. If it cannot, it propagates upward. Machine actors are never terminal.

\subsection{Bypass Mechanism: Machine Actor Exclusion}
\label{sec:bypass-mechanism}

A critical architectural property governs the escalation path: \textbf{Machine actors are excluded from the accountability chain for resolution purposes}.

\medskip
\noindent This means: an MAA node in the propagation chain \emph{may} evaluate and attempt to resolve a trigger locally (as described above), but it \emph{cannot} serve as the terminal accountability point. If the trigger reaches an MAA and that MAA cannot resolve it, the trigger does not stop — it continues propagating to the next node upward. The MAA is excluded precisely because it lacks \emph{legal authority} to authorize the class of resolution required.

\medskip
\noindent This exclusion is not a configurable parameter. It is hardwired into the protocol's architectural definition. An MAA cannot declare itself an HAA regardless of its operational authority, confidence, or the apparent resolvability of the trigger. The legal/personhood authority to bind the organization to a decision cannot be assumed by a Machine actor — it must be explicitly held by a human.

\medskip
\noindent The bypass mechanism prevents the following failure modes:

\begin{enumerate}
  \item \textbf{Compromised MAA suppression:} A maliciously compromised intermediate MAA that attempts to suppress a legitimate trigger (e.g., a sensor reading showing a Safety Envelope violation) cannot close the trigger. The trigger bypasses it and continues to the HAA.

  \item \textbf{Misconfigured MAA:} An MAA with incorrect safety parameters or an outdated model cannot erroneously resolve a trigger that should reach the Human Root. Its exclusion is automatic and structural.

  \item \textbf{Authority confusion:} An MAA that \emph{believes} it has authority to resolve a trigger — but does not — is architecturally prevented from doing so. The protocol does not rely on the MAA self-assessing its authority correctly.
\end{enumerate}

Formally, the exclusion property is expressed as:

\begin{center}
$\forall n \in \text{Nodes}: \text{IsMAA}(n) \implies \neg \text{CanBeTerminal}(n)$
\end{center}

\bigskip
\noindent The bypass rule in algorithmic form:

\begin{algorithm}
\caption{Bypass Rule — Machine Actor Exclusion}
\begin{algorithmic}[1]
\Procedure{EvaluateBypass}{$node$ $n$, TriggerPackage $pkg$}
  \If{$\text{IsHAA}(n)$}
    \State \Comment{Human Root reached — process directly}
    \State $\text{HumanRootProcess}(n, pkg)$
  \ElsIf{$\text{CanResolveLocally}(n, pkg)$}
    \State \Comment{MAA can resolve — issue directive, do NOT propagate further upward}
    \State $directive \gets \text{GenerateDirective}(n, pkg)$
    \State $\text{SendDirectiveDown}(pkg.\text{originating\_node}, directive)$
    \State $pkg.\text{status} \gets \text{RESOLVED}$
  \Else}
    \State \Comment{MAA cannot resolve — BYPASS this actor, propagate upward}
    \State $\text{PropagateUpward}(n.\text{parent}, pkg)$
  \EndIf
\EndProcedure
\end{algorithmic}
\end{algorithm}

The key property is the \textbf{else branch}: when an MAA cannot resolve a trigger, it does not absorb the trigger, log a generic error, or enter a degraded holding state. It propagates the Trigger Package upward, bypassing itself, to the next accountable node. The protocol is structurally incapable of orphaning exceptions at a Machine level.

\subsection{Timeout Handling}

The Bailout Protocol includes mandatory timeout handling to ensure that the protocol cannot enter a deadlock state in which a trigger is propagating but is never acknowledged, and to ensure that the Human Root is protected from being flooded with unprocessed triggers if multiple simultaneous conditions arise.

\medskip
\noindent\textbf{Timeout Classes:}

\begin{enumerate}
  \item[\textbf{T1 — Delivery Timeout (Originating Node $\rightarrow$ Parent):}]
  \emph{Threshold:} 30 seconds
  \emph{Trigger:} If the originating node does not receive a delivery confirmation from the parent within T1 after firing the trigger.
  \emph{Action:} Re-transmit the Trigger Package via the Local Survivability path (Section \ref{sec:local-survivability}). After 3 consecutive T1 expirations, the node escalates via all available paths simultaneously and flags the communication failure as a separate CB trigger.
\end{enumerate}

  \item[\textbf{T2 — Evaluation Timeout (Parent Node):}]
  \emph{Threshold:} 60 seconds
  \emph{Trigger:} If a parent node (MAA or HAA) does not evaluate and respond to a received Trigger Package within T2 of confirmed delivery.
  \emph{Action:} The originating node re-transmits the trigger and flags the parent node for unresponsive behavior. If the parent is an MAA and T2 expires twice consecutively, the trigger propagates to the next node in the chain above the unresponsive parent.
\end{enumerate}

  \item[\textbf{T3 — Human Acknowledgment Timeout (Human Root):}]
  \emph{Threshold:} 900 seconds (15 minutes)
  \emph{Trigger:} If the Human Root receives a Trigger Package but has not issued a directive within T3.
  \emph{Action:} A reminder alert is surfaced on the Root Dashboard. If T3 expires a second time, an SMS notification is sent to the Human Root (bypassing the dashboard entirely). T3 is the only timeout that generates an external notification, and it is limited to one SMS to prevent notification fatigue.
\end{enumerate}

  \item[\textbf{T4 — Directive Propagation Timeout:}]
  \emph{Threshold:} 120 seconds
  \emph{Trigger:} If a directive issued by the Human Root or a resolving MAA has not propagated back to the originating node within T4.
  \emph{Action:} The originating node enters a RESOLVED state and executes the directive as received, or — if the directive was not received — enters a HOLD state and re-requests the directive from the last issuer in the chain.
\end{enumerate}

\begin{algorithm}
\caption{Timeout Monitor}
\begin{algorithmic}[1]
\Procedure{StartTimeoutMonitor}{$node$ $n$, TriggerPackage $pkg$, timer\_class $T$}
  \State $deadline \gets \text{Now} + \text{Threshold}(T)$
  \State \Comment{Timeout monitor runs as independent async task}
  \State \textbf{async} \textbf{await} $\text{WaitUntil}(deadline)$

  \If{$pkg.\text{status} = \text{RESOLVED}$}
    \State \Comment{Trigger closed before timeout — monitor exits silently}
    \State \textbf{return}
  \EndIf

  \State \Switch $T$
  \Case $\text{T1}$
    \State $\text{Retransmit}(pkg.\text{parent}, pkg)$
    \State $n.\text{t1\_expiry\_count}++$
    \If{$n.\text{t1\_expiry\_count} \geq 3$}
      \State $\text{FireCommunicationFailureTrigger}(n)$
    \EndIf
  \Case $\text{T2}$
    \State $n.\text{t2\_expiry\_count}++$
    \If{$n.\text{t2\_expiry\_count} \geq 2$}
      \State $\text{BypassAndPropagate}(n.\text{parent}.\text{parent}, pkg)$
    \EndIf
  \Case $\text{T3}$}
    \State \textbf{if} $pkg.\text{t3\_reminder\_sent} = \text{false}$
    \State \Indent $\text{SurfaceReminder}(pkg)$
    \State \Indent $pkg.\text{t3\_reminder\_sent} \gets \text{true}$
    \State \textbf{else}
    \State \Indent $\text{SendSMS}($\text{Human Root}$, \text{T3 escalation alert}$)$
  \Case $\text{T4}$
    \State $n.\text{state} \gets \text{HOLD}$
    \State $\text{RequestDirectiveReIssue}(pkg.\text{directive\_issuer})$
  \EndSwitch
\EndProcedure
\end{algorithmic}
\end{algorithm}

\subsection{Local Survivability}
\label{sec:local-survivability}

In the event that communication to the parent node is disrupted — due to network failure, node crash, or malicious interference — the originating node activates its Local Survivability mode. Every node maintains a local buffer of unacknowledged Trigger Packages. When T1 expires and communication to the parent is confirmed failed, the node stores the Trigger Package locally and attempts all alternative paths to the Human Root:

\begin{enumerate}
  \item Primary path: parent uplink (Recursion Spawning)
  \item Secondary path: sibling node relay (if available and trusted)
  \item Tertiary path: direct cloud uplink (if node has independent connectivity)
  \item Quaternary path: store-and-forward (node retains trigger indefinitely until communication is restored; triggers are replayed in order upon reconnection)
\end{enumerate}

The local survivability mode ensures that the Bailout Protocol is structurally resistant to network partition. The protocol cannot be silenced by disconnecting intermediate nodes.

\section{Implementation in the BX3 Framework}
\label{sec:implementation}

The Bailout Protocol is not a safety module added to an existing autonomous system after the fact. It is the runtime expression of the BX3 Framework's upstream accountability guarantee: that every exception state at any node propagates to a human accountability anchor, bypassing all machine intermediaries, within a bounded wall-clock time. Making this guarantee architecturally enforceable rather than merely policy-level requires the Bailout Protocol to be integrated into all three functional layers at their boundaries. This section specifies that integration — how the \purpose{} anchors human accountability at the system apex, how the \bounds{} detects and fires bailout signals, how the \fact{} enforces the hard block and maintains the forensic record, and how the protocol's state machine operates as a deterministic extension of the \fact{} layer.

% ---------------------------------------------------------------
\subsection{Purpose Layer as Human Accountability Anchor}
\label{sec:purpose-anchor}

The \purpose{} layer is the only layer that may legitimately terminate a Bailout Protocol signal. This is a definitional commitment, not a design choice among alternatives. The \purpose{} is defined by its capacity for genuine accountability — the ability to bear responsibility for a decision in a way that is legally, socially, and organizationally enforceable. A machine actor, regardless of its sophistication, can optimize, recommend, and defer, but it cannot be held accountable in this sense. Therefore, only a \purpose{} actor may receive and resolve a bailout signal.

In the BX3 three-layer model, every node in a recursive system carries a \texttt{parent-pointer} and a \texttt{human-anchor} field $h(n)$ that are established at node instantiation and are immutable for the node's operational lifetime. The immutability of $h(n)$ is the structural property that terminates the escalation chain and makes the upstream accountability guarantee unconditional. No machine actor — including the node's own \bounds{}, any parent node, or any recursively spawned subordinate — may modify $h(n)$ without explicit human authorization recorded as a Forensic Ledger entry.

The \purpose{} layer serves a dual function in the Bailout Protocol. First, \textbf{authorization provenance}: the \purpose{} establishes the authoritative chain from which all legitimate system action derives. Every \bounds{} action traces back to a \purpose{} authorization. If a bailout fires, the \purpose{} owner receives the complete chain of prior authorizations and proposed actions, enabling informed judgment rather than requiring them to reconstruct context from fragments. Second, \textbf{resolution authority}: the \purpose{} owner has the authority to override the bailout, modify the node's constraints, reassign the node to a different \bounds{}, halt the system, or escalate to their own supervisory chain. No machine actor has this range of resolution options.

The key architectural property is that the \purpose{} is not merely \textit{available} as an escalation target. It is structurally the \textit{only} permissible terminal state for any unresolved exception. The BX3 upstream accountability guarantee is enforced by the layer architecture itself: the \bounds{} may not terminate a bailout signal at another machine actor, and the \fact{} hard-blocks any command that would bypass \purpose{} routing for an active exception condition. This is not a privilege level that can be reconfigured at runtime — it is a structural property of the layer model. The \purpose{} layer enforces this by preventing any node from accepting a resolution from a non-designated party when the protocol is in \texttt{HUMAN\_ACKNOWLEDGED} state.

When the Bailout Protocol fires, the \purpose{} layer performs the following operations:

\begin{enumerate}
\item Validates that $h(n)$ is defined and reachable for the originating node $n$.
\item Prevents any node in the accountability chain from re-routing the exception to an alternate human principal without an explicit \purpose{}-layer override.
\item Records the teleological context of the exception — the node's mission segment, its current objective, and the specific sub-goal that triggered the bound signal — in the authorization ledger.
\end{enumerate}

% ---------------------------------------------------------------
\subsection{Bounds Engine Triggering Bail-Out}
\label{sec:bounds-trigger}

The \bounds{} is the layer that detects and fires bailout signals. This is a direct consequence of its defining role: the \bounds{} evaluates proposed actions against encoded constraints and either approves, rejects, or escalates. When the situation exceeds its resolution authority, the \bounds{} does not guess, does not default to a safest-known action, and does not recursively attempt autonomous resolution beyond its provisioned window. It fires the Bailout Protocol.

A \bounds{} node $B$ fires a bailout signal when any of the following trigger conditions hold:

\begin{enumerate}
\item \textbf{Uncovered state:} $B$ detects a system state $s$ that is not covered by any constraint in $C_B$, where $C_B$ is the constraint set defined for $B$'s current operational scope, and $isUnresolvable(s, C_B)$ evaluates to true — meaning no rule in $C_B$ can produce an approved action for state $s$.
\item \textbf{Inter-constraint conflict:} $B$ encounters a conflict $conflict(c_1, c_2)$ between two or more constraints $c_1, c_2 \in C_B$ such that no action $a$ satisfies both constraints simultaneously.
\item \textbf{Fact Layer violation:} $B$ receives a proposed action $a$ that, when evaluated against the \fact{} validation model, would violate a deterministic enforcement rule $r \in R_{fact}$.
\item \textbf{Scope exceeded:} $B$ receives a request to perform an action $a$ that is not within the scope of $C_B$.
\item \textbf{Mandate modification request:} $B$ receives a request to modify its own constraint set $C_B$, which it cannot authorize and which therefore requires \purpose{} approval.
\end{enumerate}

These are not error codes in the conventional software engineering sense. They are \textit{boundary conditions of authority}: states in which the \bounds{} has correctly identified that it lacks the information or authorization to proceed, and that further progress requires a different kind of judgment — one that only a human accountability anchor can provide.

Critically, firing the Bailout Protocol does not indicate that the \bounds{} has made an error. The \bounds{} may be functioning correctly, following its constraints faithfully, and still encounter a condition its constraints cannot resolve. The Bailout Protocol is not a failure signal; it is an \textit{authority-exceeded} signal. It fires precisely because the \bounds{} has correctly identified the boundary of its mandate — and correctly declined to cross that boundary autonomously.

When a trigger condition fires, the \bounds{} constructs a \texttt{BailoutTrigger} object capturing the complete diagnostic context: the triggering event, the node state at time of firing, the constraint conflict or boundary condition encountered, the \bounds{} confidence level, and any proposed resolution that was attempted and failed. This object is the unit of propagation in the bailout signal chain.

The \bounds{} also enforces an atomic suspension property: upon firing, the originating agent is atomically suspended from further action on the affected task. This prevents a second agent or process from racing past the bailout to execute the same action on a different thread while the escalation is in progress. The suspension is enforced by the \fact{} upon receipt of the \texttt{BailoutTrigger} — not by the \bounds{} itself — ensuring that the suspension cannot be suppressed by a compromised \bounds{}.

The interaction between \bounds{} and the protocol follows a strict temporal layering:

\begin{enumerate}
\item The \bounds{} layer detects and classifies the bound signal, constructing the \texttt{BailoutTrigger}.
\item The \fact{} applies the hard block to the originating node's physical actuator commands.
\item The \bounds{} attempts autonomous resolution within $T_{\text{bounds}}$ (default 60 seconds).
\item The \fact{} evaluates each proposed resolution against the forensic ledger and issues an \texttt{approved} or \texttt{denied} verdict.
\item If resolution fails or $T_{\text{bounds}}$ expires, the Bailout Protocol receives a \texttt{bailout-triggered} signal and transitions to \texttt{ESCALATING}.
\item The protocol propagates the exception to $h(n)$ via the optimal pathway.
\end{enumerate}

The \bounds{} layer's role terminates at step 5. It does not propagate exceptions — that is the protocol's responsibility. It does not select communication pathways — the \fact{} layer's \texttt{SelectPathway} function performs that operation. And it does not wait for human acknowledgment — the protocol manages the \texttt{HUMAN\_ACKNOWLEDGED} state independently. This separation of concerns is deliberate: the \bounds{} defines the policy of what constitutes a violation and how to resolve it; the Bailout Protocol enforces the procedure by which unresolved violations reach the human principal.

% ---------------------------------------------------------------
\subsection{Fact Layer Enforcing Audit Trail}
\label{sec:fact-enforcement}

The \fact{} layer serves two distinct functions in the Bailout Protocol: it provides the deterministic enforcement mechanism that prevents premature or unauthorized continuation of machine action during a bailout, and it maintains the Forensic Ledger that records the complete bailout event chain for post-hoc accountability review.

The \fact{} enforces a \textit{hard block} on any physical actuator command associated with a live bailout signal. When a \texttt{BailoutTrigger} is in the \texttt{fired} or \texttt{escalating} state, the \fact{} refuses to forward any \bounds{} command to physical actuators. This is not a conditional software check:

\begin{verbatim}
if (bailoutActive) { block() }  // SOFT BLOCK -- bypassable
\end{verbatim}

It is a deterministic enforcement rule in the \fact{} state machine:

\begin{verbatim}
on receiving cmd to actuator:
    if bailout_state != RESOLVED:
        HARD_BLOCK  // structurally enforced
        return
    else:
        forward cmd
\end{verbatim}

The block cannot be overridden by a higher-confidence \bounds{} proposal, a changed business context, a timeout, or any machine actor. The block holds until the \texttt{BailoutTrigger} reaches \texttt{human\_review} and the \purpose{} owner issues an explicit resolution. The block is a property of the \fact{}'s state machine, not a conditional check that can be bypassed by a sufficiently confident proposal.

The Forensic Ledger records every bailout event in real time across three planes of the nine-plane DAP structure (Plane 4: Accountability; Plane 6: Temporal Sequencing; Plane 8: Compliance Posture). A dedicated \texttt{BailoutEventRecord} schema extends the standard event record with the following fields:

\begin{itemize}
\item \textbf{Entry type:} A machine-readable code identifying the category of trigger condition (e.g., \texttt{UNCOVERED\_STATE}, \texttt{INTERCONSTRAINT\_CONFLICT}, \texttt{FACT\_LAYER\_VIOLATION}, \texttt{SCOPE\_EXCEEDED}).
\item \textbf{Originating node:} The cryptographic identifier of the \bounds{} node at which the bailout was fired, enabling reconstruction of the full propagation path.
\item \textbf{Diagnostic snapshot:} A complete capture of the \bounds{} state at time of firing, including constraint state, sensor readings, prior actions, and the specific conflict or boundary condition encountered.
\item \textbf{Propagation chain:} An ordered, timestamped list of every node through which the \texttt{BailoutTrigger} propagated, with the action taken at each hop.
\item \textbf{Resolution record:} The resolution issued by the \purpose{} owner, including the action taken, any mandate modifications issued, and the explicit reasoning provided.
\item \textbf{Chain integrity hash:} The SHA-256 hash of this ledger entry, chained to the hash of the preceding entry, ensuring tamper evidence across the full bailout record.
\end{itemize}

The append-only, hash-chained structure ensures that any modification to a prior entry — whether accidental or deliberate — breaks the hash chain at the point of modification and is detectable in $O(1)$ by any verifier performing a chain integrity check. The \fact{} performs this check automatically after every event write and triggers a bailout if a break is detected, ensuring that ledger integrity is itself subject to the upstream accountability guarantee.

The \fact{} layer additionally maintains the Pathway Health Registry (PHR), a runtime data structure that tracks the availability and latency of each communication pathway to the human anchor. Before the protocol propagates an exception, it queries the \fact{} via \texttt{SelectPathway}, which returns the pathway with the lowest estimated latency among currently functional pathways. The \fact{} updates the PHR continuously via background health-ping threads, ensuring that pathway selection is informed by current availability rather than stale state.

The following functions are provided exclusively by the \fact{} layer within the Bailout Protocol:

\begin{enumerate}
\item \texttt{LedgerWrite(entry)} — appends an authorization ledger entry for each state transition and resolution event.
\item \texttt{LedgerVerify(node, action)} — evaluates whether an action falls within the node's ledger-attested permission scope.
\item \texttt{SelectPathway(node)} — queries the Pathway Health Registry and returns the optimal functional pathway to $h(n)$.
\item \texttt{SealLedger()} — applies the cryptographic sealing operation that makes the ledger tamper-evident.
\end{enumerate}

% ---------------------------------------------------------------
\subsection{Bailout Protocol as Runtime Expression of BX3's Upstream Accountability Guarantee}
\label{sec:runtime-expression}

The upstream accountability guarantee — that systems built on the BX3 Framework \textit{fail upward into human consciousness, never downward into algorithmic chaos} — is a structural property of the framework, guaranteed by the layer architecture. The Bailout Protocol is the runtime mechanism by which this guarantee is effectuated in real-time execution.

Consider the contrast with conventional autonomous systems. In a conventional system, human oversight is typically a supervisory layer that monitors agent outputs and can intervene upon detecting a problem. The pipeline continues until a threshold is crossed. This creates two failure modes that BX3 eliminates.

First, the system can fail dangerously while awaiting threshold crossing if the threshold is miscalibrated or the failure mode was unanticipated. Second, human oversight in the conventional model is reactive rather than mandatory: it depends on a human noticing a problem in time to intervene, with all the attentional and cognitive limitations that entails.

The BX3 upstream accountability guarantee eliminates both failure modes by making human oversight the \textit{mandatory terminal state of every unresolvable condition}, not an optional intervention layer. The Bailout Protocol fires when the \bounds{} identifies a condition it cannot resolve — not when a human notices a problem. The human is inserted into the loop by the architecture, not by their own vigilance. This transforms the human's role from reactive monitor to active, authoritative decision-maker engaged precisely at the moments when the system's autonomous capability has reached its boundary.

The protocol does not merely \textit{enable} upstream accountability. It \textit{enforces} it architecturally. The distinction is between a system designed to route exceptions to a human \textit{if possible}, and a system architecturally incapable of terminating an unresolvable exception at any machine actor. The BX3 Framework is the latter. The Bailout Protocol is the mechanism that makes this architectural commitment operative at runtime.

The relationship is asymmetric but essential. The protocol does not generate accountability — it propagates it. The \purpose{} layer owns accountability; the protocol is the instrument through which accountability is exercised. The \fact{} layer independently certifies it. Together, the three layers and the protocol constitute a closed-loop system in which the accountability guarantee is (a) architecturally declared by the \purpose{} layer, (b) enforced at runtime by the protocol, and (c) independently verified and recorded by the \fact{} layer.

The protocol also resolves the non-terminating delegation problem that arises in deep agent networks. Without the protocol, a bound signal at depth $d$ in a network of $N$ nodes with $N \gg d$ could propagate indefinitely upward through machine actors, each of which might attempt its own resolution and re-propagate. The protocol terminates this cycle by enforcing a strict resolution hierarchy: the \bounds{} layer attempts resolution for $T_{\text{bounds}}$, the protocol propagates upward on failure, and at \texttt{HUMAN\_ACKNOWLEDGED} state no machine actor may issue or substitute a resolution. The human principal is the terminal resolution authority by architectural design, not by fallback.

% ---------------------------------------------------------------
\subsection{Forensic Ledger Recording}
\label{sec:forensic-ledger}

Every state transition in the Bailout Protocol generates one or more authorization ledger entries in the \fact{} layer. These entries are the forensic record of the protocol's execution and form the evidentiary basis for post-hoc audit, compliance verification, and regulatory review.

Each entry contains the following fields:

\begin{enumerate}
\item \textbf{Entry type}: the protocol event that generated the entry (e.g., \texttt{BOUNDED}, \texttt{ESCALATING}, \texttt{HUMAN\_ACKNOWLEDGED}, \texttt{RESOLVED}, \texttt{RESOLUTION\_DENIED}).
\item \textbf{Node identity}: the cryptographic identifier of the node at which the event occurred.
\item \textbf{Timestamp}: wall-clock time of the event, sourced from a synchronized time service.
\item \textbf{Trigger context}: the bound signal's classification, the input that triggered it, and the node's current operational state at the time of firing.
\item \textbf{Propagation path}: the ordered list of nodes through which the exception propagated en route to $h(n)$.
\item \textbf{Pathway selection}: the pathway selected by \texttt{SelectPathway} and its PHR-assessed latency at the time of selection.
\item \textbf{Human directive}: the directive issued by $h(n)$ in \texttt{HUMAN\_ACKNOWLEDGED} state, if applicable, including the directive type (\texttt{ACK}, \texttt{RESOLVE}, \texttt{REJECT}) and any binding parameters.
\item \textbf{Resolution outcome}: whether the exception was resolved, acknowledged, or left unresolved, and the \fact{} attestation of the resolution's ledger conformance.
\end{enumerate}

The ledger records not only successful resolutions but also failures: denials by the \fact{} layer, timeouts at any layer, pathway degradation events, and human directives that override the protocol's recommended resolution. This completeness is critical for auditability — a forensic record that only shows approvals is not a reliable account of system behavior; it is a curated selection.

The ledger's append-only property is enforced by the cryptographic chaining mechanism. Each entry's digest incorporates the SHA-256 digest of the preceding entry, forming a chain in which any entry's historical position can be verified by recomputing the hash chain from the genesis entry to the target entry. The \fact{} provides \texttt{VerifyChain(start, end)} — a function that recomputes digests over the specified range and returns a boolean indicating whether the chain is intact. This verification can be performed by any external auditor without access to the \fact{}'s internal state, using only the publicly accessible hash values.

For regulatory contexts in which the evidentiary chain must be established to a specific standard of certainty, the \fact{} layer additionally supports attestation notarization: a human principal may invoke a \texttt{NotarizeLedger(from, to)} operation that generates a cryptographically signed attestation of the chain's integrity over the specified entry range, suitable for submission as evidence in compliance proceedings.

% ---------------------------------------------------------------
\subsection{State Machine as Fact Layer Extension}
\label{sec:state-machine}

The Bailout Protocol state machine $\mathcal{B}$ is hard-encoded in the \fact{} layer. The \fact{} does not merely enforce the output of the Bailout Protocol; it defines the protocol's state machine itself. This is a critical architectural commitment: the state transitions of the Bailout Protocol are deterministic, not probabilistic, and are enforced by the same mechanism that enforces all \fact{} constraints. There is no separate Bailout Protocol module that can be modified independently of the \fact{}.

The state machine defines the following states:

\begin{enumerate}
\item \textbf{\texttt{idle}}: Default state. No trigger is active. The \bounds{} operates normally within its constraint set. No ledger entry required.
\item \textbf{\texttt{fired}}: A trigger condition has been detected. The \bounds{} has constructed the \texttt{BailoutTrigger} and the \fact{} has received it. A \texttt{BOUNDED} ledger entry is created, recording the trigger context and the start of $T_{\text{bounds}}$. All physical actuator commands from the affected node are hard-blocked.
\item \textbf{\texttt{escalating}}: The \texttt{BailoutTrigger} has been dispatched to the parent node in the propagation chain. A \texttt{ESCALATING} ledger entry is created for each hop, recording the pathway selected, the transmit timestamp, and the acknowledgment status. The \fact{} maintains the hard block on the originating node.
\item \textbf{\texttt{human\_review}}: The \texttt{BailoutTrigger} has arrived at the \purpose{} human anchor $h(n)$. The \purpose{} owner has the diagnostic context and is actively reviewing the bailout. A \texttt{HUMAN\_ACKNOWLEDGED} ledger entry is created, recording the acknowledgment timestamp and the identity of the human principal who acknowledged. The hard block remains in effect.
\item \textbf{\texttt{resolved}}: The \purpose{} owner has issued a resolution. A \texttt{RESOLVED} ledger entry is created, recording the directive type and any binding parameters. The \fact{} records the resolution in the Forensic Ledger, releases the hard block, and the system transitions back to \texttt{idle}.
\end{enumerate}

Each state in $\mathcal{B}$ is therefore not merely a protocol condition — it is a forensic marker in the \fact{} layer's authorization ledger. The state machine provides the transition logic; the \fact{} layer provides the evidentiary record. The separation is intentional: the protocol's correctness is verified by the state machine's transition rules, while the evidence of its execution is preserved independently in the ledger.

The state transition function is:

\begin{equation}
\delta: \mathcal{S} \times \mathcal{E} \rightarrow \mathcal{S}
\end{equation}

where $\mathcal{S} = \{idle, fired, escalating, human\_review, resolved\}$ and $\mathcal{E}$ is the set of events defined by the transition predicates in Table~\ref{tab:bailout-transitions}. Each transition is atomic and deterministic: given the same state and event, the same next state is produced.

\begin{table}[H]
\centering
\caption{Bailout Protocol State Transitions as Deterministic \fact{} Rules}
\label{tab:bailout-transitions}
\begin{tabular}{@{}lp{5.5cm}l@{}}
\toprule
\textbf{From} & \textbf{Event (Transition Predicate)} & \textbf{To} \\
\midrule
$idle$ & $trigger\_condition() = \top$ & $fired$ \\
$fired$ & Diagnostic context construction complete & $escalating$ \\
$fired$ & $T_{\text{bounds}}$ timeout expires & $escalating$ \\
$escalating$ & Parent node is \purpose{} anchor, trigger received & $human\_review$ \\
$escalating$ & Parent node is not \purpose{} anchor & $escalating$ (re-dispatch) \\
$human\_review$ & \purpose{} owner issues resolution & $resolved$ \\
$human\_review$ & $T_{\text{human}}$ timeout expires & $escalating$ (re-notify) \\
$resolved$ & Post-processing complete & $idle$ \\
\bottomrule
\end{tabular}
\end{table}

The transition $human\_review \rightarrow resolved$ requires a \purpose{} authorization payload: a resolution authored by a named human and digitally signed with their \purpose{} credentials. The \fact{} verifies this payload before accepting the transition. No machine actor can produce a \purpose{} authorization payload, so no machine actor can resolve a bailout unilaterally.

The three state invariants maintained by the \fact{}'s enforcement architecture are:

\begin{align}
\mathrm{INV}_1 &: \forall n \in \mathrm{nodes},\; \mathrm{state}(n) \in \mathcal{S} \setminus \{\mathrm{escalating}\} \iff \neg\mathrm{active\_bailout}(n). \tag{INV-1}
\end{align}

INV-1 is enforced by the \fact{}'s pre-commit hook: a node may enter a non-escalating state only when the Bailout Monitor confirms that no active propagation is in flight. The state and the active-bailout flag are updated atomically, preserving the equivalence.

\begin{align}
\mathrm{INV}_2 &: \mathrm{state}(n) = \mathrm{fired} \implies \mathrm{active\_bailout}(n) \land \delta(\mathrm{fired}, e_{\mathrm{timeout}}) = \mathrm{escalating}. \tag{INV-2}
\end{align}

INV-2 follows from the \fact{}'s enforcement of the transition relation. Any node in \texttt{fired} has an active bailout, and the only terminal transition defined for this state with an enabled event is to \texttt{escalating}. The \fact{} enforces this transition compulsorily upon receipt of $e_{\mathrm{timeout}}$; no machine actor may interpose.

\begin{align}
\mathrm{INV}_3 &: \mathrm{state}(n) = \mathrm{human\_review} \implies \exists\; \mathrm{Purpose\_layer\_auth}(n) \land \mathrm{directive}(n) \in \{\texttt{ACK}, \texttt{RESOLVE}, \texttt{REJECT}\}. \tag{INV-3}
\end{align}

INV-3 establishes that a node enters \texttt{human\_review} only when a \purpose{}-layer-authorized directive is present and recorded. This prevents the state machine from entering the human-review state based on machine-generated signals alone; the state is only reached upon receipt of an authenticated \purpose{} communication.

The hard block is monotonic with respect to the state machine: the block is applied at \texttt{fired} and is not released until \texttt{resolved}. There is no state in which the hard block is conditionally applied based on confidence level, elapsed time, or business context. This monotonicity property is what makes the Bailout Protocol a structurally reliable safety mechanism rather than a configurable policy that can be inadvertently weakened.

The \fact{} maintains the state machine for every active \texttt{BailoutTrigger} across the entire system topology, enabling the framework to support recursive spawning patterns in which child nodes inherit the Bailout Protocol state machine from their parent, with the same states, transitions, and enforcement rules. The protocol is not a single-instance mechanism — it is a replicated, per-node state machine that propagates upward through the hierarchy when triggered.

The combination of the state machine specification in the \fact{}, the monotonic hard block on actuator commands, the three-plane Forensic Ledger recording, the \bounds{} trigger conditions, and the three state invariants constitutes a complete, architecturally enforced implementation of the Bailout Protocol. The upstream accountability guarantee is not maintained by convention, by operational policy, or by the reliability of any individual component. It is maintained by the layer architecture itself, enforced at the \fact{}, and triggered by the \bounds{}. The system fails upward, without exception, and with a complete forensic record of every step in the chain.

\subsection{Summary of Structural Guarantees}

The Bailout Protocol provides three interlocking structural guarantees that no prior autonomous system has achieved:

\medskip
\noindent\textbf{1. Indefeasible Human Accountability:} The Human Root is the sole terminal accountability point. No Machine actor can close a trigger that requires human judgment. This property is hardwired and cannot be overridden by any configuration change, system update, or compromise of intermediate nodes.

\medskip
\noindent\textbf{2. No Silent Failures:} Every condition a node cannot resolve with certainty generates a Trigger Package. There is no condition category for which the node is permitted to silently absorb the exception and continue. Either the node resolves it locally (with full logging), or it propagates it upward.

\medskip
\noindent\textbf{3. Partition Resistance:} The Local Survivability mode ensures that trigger propagation is not disrupted by network partition. Triggers accumulate locally and are replayed upon reconnection. A partitioned node cannot accumulate unresolved exceptions indefinitely without the Human Root being aware once connectivity is restored.

These guarantees make the BX3 architecture legally defensible in high-stakes domains — agriculture, healthcare, infrastructure, finance — where autonomous action without accountable human authorization creates unacceptable legal and societal risk. The Bailout Protocol is not a safety feature. It is a structural constraint that makes human accountability architecturally indestructible.

\medskip
\noindent\textbf{Next:} Section \ref{sec:forensic-ledger} establishes the cryptographic audit infrastructure that records every state transition, trigger firing, propagation step, directive issued, and physical outcome across the three-plane Ledger system.

\end{document}
